% Baccalauréat malgache, série D, 2017 : énoncé.
\TCDBacAttribution{D}{2017}
\providecommand{\Plan}{(\mathscr{P})}
\providecommand{\Ouv}{(O ; \vec{u}, \vec{v})}
\providecommand{\cm}{~\mathrm{cm}}
\providecommand{\Courbe}{(\mathscr{C})}
\providecommand{\CourbeP}{(\mathscr{C}')}
\providecommand{\D}{(\mathscr{D})}
\providecommand{\dx}{~\mathrm{d}x}

\noindent \underline{\textsc{Exercice I}} (5 points) 
      \begin{enumerate}
            \item \begin{enumerate}
                        \item R\'esoudre dans $\mathbb{C}$, l'\'equation 
                              \begin{equation}
                                    \TCDtag{$E_1$}
                                          z^2 - 2 z + 2 = 0      \;.
                                    \label{eq 1: Exo 1 Bac D 2017}
                              \end{equation}
                        \item Pr\'eciser le module et un argument de chacune des solutions de \eqref{eq 1: Exo 1 Bac D 2017}.
                        \item En d\'eduire les solutions dans $\mathbb{C}$ de l'\'equation 
                              \begin{equation}
                                    \TCDtag{$E_2$}
                                          (-iz+3i+3)^2 - 2(-iz+3i+3) + 2 = 0
                                    \label{eq 2: Exo 1 Bac D 2017}
                              \end{equation}
                              (on pourra poser $Z = -iz+3i+3$).
                  \end{enumerate}
            \item Le plan $\Plan$ est rapport\'e \`a un rep\`ere orthonorm\'e direct $\Ouv$ d'unit\'e $1\cm$.
                  On donne les points $A$, $B$ et $C$ d'affixes respectives : 
                  $z_A = 1 + i$, $z_B = 1 - i$ et $z_C = 2 - 2i$.
                  \begin{enumerate}
                        \item Placer les points $A$, $B$ et $C$.
                        \item Donner la forme trigonom\'etrique de $U = \frac{z_E - z_C}{z_E - z_A}$
                              avec $E$ d'affixe $z_E = 3$.
                        \item En d\'eduire la nature du triangle $AEC$.
                  \end{enumerate}
            \item On consid\`ere la transformation $R$ du plan $\Plan$ dans $\Plan$ qui \`a tout point $M$
                  d'affixe $z = x + iy$ associe le point $M'$ d'affixe $z' = x' + i y'$ tel que : 
                        $$ \begin{multiform} x' = -y \\ y' = x - 4 \;. \end{multiform} $$
                  \begin{enumerate}
                        \item Donner l'expression complexe de $R$.
                        \item Pr\'eciser la nature et les \'el\'ements caract\'eristiques de $R$.
                        \item Soit $(\Gamma)$ le cercle de centre $E$ et de rayon $\sqrt{5}$.
                              D\'eterminer et construire l'image $(\Gamma ')$ de $(\Gamma)$ par $R$ 
                              dans le rep\`ere pr\'ec\'edent.
                  \end{enumerate}
      \end{enumerate}
%%%%%%%%%%%%%%%%%%%%%%%%%%%%%%%%%%%%%%%%%%%%%%%%%%%%%%%%%%%%%%%%%%%%%%%%%%%%%%%%%%%%%%%%%%%%%%%
%%%%%%%%%%%%%%%%%%%                   %%%%%%%%%%%%%%%%%%%                   %%%%%%%%%%%%%%%%%%%
%%%%%%%%%%%%%%%%%%%%%%%%%%%%%%%%%%%%%%%%%%%%%%%%%%%%%%%%%%%%%%%%%%%%%%%%%%%%%%%%%%%%%%%%%%%%%%%
\noindent \underline{\textsc{Exercice II}} (5 points) 
      \begin{description}
            \item[A. ] On dispose deux urnes dont l'une $U_1$ contient 8 boules num\'erot\'ees : 
                  0; 2; 2; 2; 3; 3; 4; 5 
                  et l'autre $U_2$ contient 5 boules num\'erot\'ees : 0; 0; 1; 2; 2.
                  Les boules sont indiscernables au toucher.
                  \begin{enumerate}
                        \item On tire au hasard et successivement sans remise 3 boules de l'urne $U_1$.
                              Calculer la probabilit\'e de chacun des \'ev\`enements suivants :
                              \begin{itemize}
                                    \item $E_1$ : << Les num\'eros des 3 boules tir\'ees sont pairs >>.
                                    \item $E_2$ : << Avoir 3 boules dont la somme des num\'eros tir\'es est \'egale \`a 6 >>.
                              \end{itemize}
                        \item On tire au hasard et simultan\'ement 2 boules de $U_2$ et 1 boule de $U_1$.
                              On suppose que les \'ev\`enements \'el\'ementaires sont \'equiprobables.
                              Soit $X$ la variable al\'eatoire \'egale au nombre de boules portant le num\'ero 2.
                              \begin{enumerate}
                                    \item Donner la loi de probabilit\'e de $X$.
                                    \item Calculer l'\'esp\'erance math\'ematique $E(X)$ de $X$.
                              \end{enumerate}
                  \end{enumerate}
            \item[B. ] Lors d'une \'epreuve facultative d'EPS, on donne pour 4 \'el\`eves les points bonus $x_i$ 
                  obtenus en \'epreuves de ``course de fond'' et $y_i$ obtenus en \'epreuve de ``saut''.

                        \begin{tabular}{|c|c|c|c|c|}
                              \hline
                              $x_i$ & 1 & 2 & 5 & 7 \\
                              \hline
                              $y_i$ & 2 & 2 & 4 & 5 \\
                              \hline
                        \end{tabular}
                  \begin{enumerate}
                        \item Calculer le coefficient de corr\'elation lin\'eaire. Interpr\'eter le r\'esultat.
                        \item D\'eterminer, par la m\'ethode des moindres carr\'es l'\'equation de la droite
                              de r\'egression de $y$ en $x$.
                        \item Estimer le bonus en \'epreuve de saut pour un \'el\`eve ayant 9 points de bonus en course de fond.
                  \end{enumerate}
      \end{description}
%%%%%%%%%%%%%%%%%%%%%%%%%%%%%%%%%%%%%%%%%%%%%%%%%%%%%%%%%%%%%%%%%%%%%%%%%%%%%%%%%%%%%%%%%%%%%%%
%%%%%%%%%%%%%%%%%%%                   %%%%%%%%%%%%%%%%%%%                   %%%%%%%%%%%%%%%%%%%
%%%%%%%%%%%%%%%%%%%%%%%%%%%%%%%%%%%%%%%%%%%%%%%%%%%%%%%%%%%%%%%%%%%%%%%%%%%%%%%%%%%%%%%%%%%%%%%
\noindent \underline{\textsc{Probleme}} \newline
      On consid\`ere la fonction $f$ d\'efinie sur $\mathbb{R}$ par : 
      $ \begin{multiform}
            f(x) = 1 - \expo{x}   \text{ si } x \leq 0 \\ 
            f(x) = x (-1+\ln x)   \text{ si } x > 0    \;.
      \end{multiform} $ \newline
      On note par $\Courbe$ sa courbe repr\'esentative dans un plan muni d'un rep\`ere orthonorm\'e \Oij{} d'unit\'e $1\cm$.
      \begin{enumerate}
            \item Etudier la continuit\'e de $f$ en $x_0 = 0$.
            \item \begin{enumerate}
                        \item Montrer que 
                              $\displaystyle{\lim_{x\rightarrow 0^-} \left[ \frac{f(x)-f(0)}{x} \right] = -1}$ et 
                              $\displaystyle{\lim_{x\rightarrow 0^+} \left[ \frac{f(x)-f(0)}{x} \right] = -\infty}$.
                        \item Interpr\'eter graphiquement ces r\'esultats.
                        \item Calculer les limites de $f$ en $-\infty$ et $+\infty$.
                  \end{enumerate}
            \item \begin{enumerate}
                        \item Pour $x>0$, calculer $f'(x)$ et \'etudier son signe.
                        \item Pour $x\leq 0$, calculer $f'(x)$ et \'etudier son signe.
                        \item Dresser le tableau de variation de $f$.
                  \end{enumerate}
            \item Donner l'\'equation de la tangente $(T)$ \`a $\Courbe$ au point d'abscisse $x_0 = e$.
            \item \begin{enumerate}
                        \item Calculer $f(e^2)$.
                        \item Etudier les branches infinies de $\Courbe$.
                  \end{enumerate}
            \item Tracer la tangente $(T)$ et la courbe $\Courbe$ en pr\'ecisant les demi-tangentes au point $O$.
            \item Soit $g$ la restriction de $f$ sur l'intervalle $I = [ e ; +\infty [$.
                  \begin{enumerate}
                        \item Montrer que $g$ admet une fonction r\'eciproque $g^{-1}$ d\'efinie sur l'intervalle $J$
                              que l'on pr\'ecisera.
                        \item Dresser le tableau de variation de $g^{-1}$.
                        \item Tracer la courbe $\CourbeP$ de $g^{-1}$ dans le m\^eme rep\`ere que $\Courbe$.
                        \item Calculer $(g^{-1})' (e^2)$.
                  \end{enumerate}
            \item En utilisant une int\'egration par parties, calculer, en $\cm^2$, l'aire $\mathcal{A}$ du domaine plan
	                  limit\'e par la courbe $\Courbe$, l'axe des abscisses et les droites d'\'equations $x=e$ et $x=\expo{2}$.
                  
                  On donne $e \simeq 2,7$ et $\expo{2} \simeq 7,4$.
      \end{enumerate}
      









